% Unit 088 terjemahan Bahasa Indonesia dari chapter7.tex baris 1--418.
% Sumber: Methods of Algebra, Volume 2, commit
% 9a5803ff77dd3257484cb177f851a73770a59dd3 (CC BY 4.0).
% stable-unit-id: o014.aljabr2.chapter7.monoidal-algebras
% Cakupan: pembukaan Bab 7 dan Bagian 7.1 lengkap tentang aljabar dalam
% kategori monoidal; berhenti tepat sebelum Bagian 7.2 pada baris 419.

% segment-id: o014.li2.chapter7.monoidal-algebras.h001
\chapter{Teori Monad}\label{sec:monads}
\hypertarget{o014.aljabr2.chapter7.monoidal-algebras}{ }

% segment-id: o014.li2.chapter7.monoidal-algebras.p001
Titik awal bab ini, \S\ref{sec:alg-in-monoidal-cat}, ialah aljabar di atas
kategori monoidal $(\mathcal{V},\otimes)$. Yang dimaksud ialah objek
$A\in\Obj(\mathcal{V})$ yang dilengkapi perkalian
$\mu:A\otimes A\to A$ dan morfisme unit $\eta:\mathbf{1}\to A$, serta
	memenuhi hukum asosiativitas dan hukum unit. Konsep ini merupakan padanan konsep monoid,
sedangkan aljabar-$\Bbbk$ di atas gelanggang komutatif $\Bbbk$ merupakan
contoh tipikal. Aljabar-aljabar tersebut membentuk kategori
$\cate{Alg}(\mathcal{V})$. Bila $(\mathcal{V},\otimes)$ memiliki sifat
komutatif yang lebih kuat, $\cate{Alg}(\mathcal{V})$ pun memiliki struktur dan
sifat tambahan. Situasi terbaik terjadi ketika $\mathcal{V}$ merupakan
kategori monoidal simetris; dalam hal ini $\cate{Alg}(\mathcal{V})$ juga
merupakan kategori monoidal simetris.

% segment-id: o014.li2.chapter7.monoidal-algebras.p002
Untuk suatu aljabar $A$, modul-$A$ kiri ialah objek $M$ dari $\mathcal{V}$
bersama morfisme perkalian skalar $\mu_M:A\otimes M\to M$ yang memenuhi hukum
asosiativitas dan unit; modul kanan dan bimodul didefinisikan serupa. Karena
	kompleks menjadi objek utama buku ini, langkah yang wajar ialah
mengambil kategori abelian $\mathcal{A}$ dengan koproduk terhitung serta
bifunktor aditif eksak kanan
$\otimes:\mathcal{A}\times\mathcal{A}\to\mathcal{A}$, memperluas $\otimes$
ke tingkat kompleks dengan kompleks total, lalu meninjau kategori monoidal
$(\cate{C}(\mathcal{A}),\otimes)$ yang biasanya diberi struktur anyaman Koszul
\eqref{eqn:Koszul-braiding}. Dengan demikian, aljabar-dg dan modul-dg muncul
secara alami sebagai kelanjutan teori; ``dg'' adalah singkatan dari
``diferensial bergradasi''. Inilah pokok bahasan
\S\ref{sec:dga-monoidal}.

% segment-id: o014.li2.chapter7.monoidal-algebras.p003
Kita juga mengetahui bahwa himpunan $\Hom$ dalam $\cate{C}(\mathcal{A})$ dapat
diperluas menjadi kompleks $\Hom$, dan komposisi morfisme dapat diangkat
ke tingkat kompleks $\Hom$; khususnya, kompleks endomorfisme
$\End^\bullet(X)$ suatu objek $X$ menjadi aljabar-dg. Generalisasi pengamatan
ini menghasilkan gagasan kategori-dg (Definisi~\ref{def:dgCat}), dan antara
kategori-dg dapat dibicarakan funktor-dg. Setelah pengantar singkat mengenai
kategori monoidal tertutup dalam \S\ref{sec:closed-monoidal}, pada
\S\ref{sec:dg-closed} akan dijelaskan bahwa kategori $\cate{dgCat}$ dari semua
kategori-dg kecil bersifat tertutup. Artinya, seluruh funktor-dg antara dua
kategori-dg kecil $\mathcal{C}$ dan $\mathcal{D}$ membentuk suatu
kategori-dg $\iHom(\mathcal{C},\mathcal{D})$.

% segment-id: o014.li2.chapter7.monoidal-algebras.p004
Berbagai struktur-dg muncul secara alami sekaligus menjadi alat penting dalam
aljabar linear. Namun, bab ini hanya memberikan pengantar ringkas, dan
\S\ref{sec:bialgebra} akan beralih kepada koaljabar dan komodul di atas
kategori monoidal $\mathcal{V}$. Keduanya tidak lain adalah aljabar dan modul
di atas kategori monoidal lawan $\mathcal{V}^{\opp}$. Secara lebih konkret,
koaljabar $C$ dilengkapi koperkalian $\Delta:C\to C\otimes C$ dan morfisme
kounit $\epsilon:C\to\munit$, sedangkan komodul kiri $M$ dilengkapi
koperkalian skalar $\rho:M\to C\otimes M$; hukum asosiativitas dan unitnya
merupakan versi dengan semua panah dibalik.

% segment-id: o014.li2.chapter7.monoidal-algebras.p005
Jika $\mathcal{V}$ merupakan kategori monoidal beranyam, struktur aljabar dan
koaljabar yang saling kompatibel $(A,\mu,\eta,\Delta,\epsilon)$ pada objek $A$
disebut bialjabar, dan bialjabar yang dilengkapi antipoda disebut aljabar Hopf
(Definisi~\ref{def:Hopf-algebra}). Dalam kasus konkret
$\mathcal{V}=\Bbbk\dcate{Mod}$, konsep-konsep ini diilhami oleh geometri dan
topologi. Sebagai contoh, kohomologi grup Lie, atau lebih umum grup-H, secara
alami menjadi aljabar Hopf; kelompok kuantum merupakan kelas contoh penting
lainnya. Selain itu, aljabar grup $\Bbbk[G]$ dari sembarang grup $G$ juga
merupakan aljabar Hopf.

% segment-id: o014.li2.chapter7.monoidal-algebras.p006
Sekarang kita memperkenalkan pelaku utama \S\ref{sec:Beck}, yang sekaligus
menjadi judul bab ini: monad. Untuk sembarang kategori $\mathcal{C}$, kategori
endofunktor $\mathcal{C}^{\mathcal{C}}$ menjadi kategori monoidal ketat terhadap
komposisi, dan monad tidak lain adalah aljabar dalam
$\mathcal{C}^{\mathcal{C}}$, sedangkan komonad merupakan versi dualnya.
Walaupun sederhana, konsep ini memainkan peranan penting dalam aljabar,
geometri, dan ilmu komputer. Salah satu kegunaan menonjol monad ialah
menjelaskan berbagai ``resolusi standar'' dalam aljabar homologis, misalnya
$\mathsf{B}R$ yang muncul dalam \S\ref{sec:HH} atau $\mathsf{L}$ dalam
\S\ref{sec:G-mod}; aspek ini akan dikaji dalam
\S\ref{sec:bar-resolution}.

% segment-id: o014.li2.chapter7.monoidal-algebras.p007
Monad dan komonad terutama berasal dari pasangan adjungsi
% segment-id: o014.li2.chapter7.monoidal-algebras.g001
\[\begin{tikzcd}
	F: \mathcal{C} \arrow[shift left, r] & \mathcal{D}: G \arrow[shift left, l].
\end{tikzcd}\]
Pasangan ini menentukan monad $T$ pada $\mathcal{C}$ dan komonad $L$ pada
$\mathcal{D}$, yang melalui unit $\eta$ dan kounit $\varepsilon$ dari adjungsi
dinyatakan sebagai
% segment-id: o014.li2.chapter7.monoidal-algebras.g002
\[\begin{array}{|c|c|}\hline
	\text{monad} & \text{komonad} \\ \hline
	T := GF & L := FG \\
	\mu := \left[ T^2 \xrightarrow{G\varepsilon F} T \right] & \delta := \left[ L \xrightarrow{F\eta G} L^2 \right] \\
	\eta : \identity_{\mathcal{C}} \to T & \varepsilon: L \to \identity_{\mathcal{D}} \\ \hline 
\end{array}\]
Pasangan adjungsi tersebut secara alami menginduksi funktor
$\mathbb{K}:\mathcal{D}\to\mathcal{C}^T$ (atau
$\mathcal{C}\to\mathcal{D}^L$), dengan $\mathcal{C}^T$ (atau
$\mathcal{D}^L$) menyatakan kategori modul di bawah aksi monad (atau versi
komonadnya). Jika $\mathbb{K}$ merupakan ekuivalensi, pasangan adjungsi itu
disebut monadik (atau komonadik). Ini berarti informasi dalam $\mathcal{D}$
(atau $\mathcal{C}$) dapat direkonstruksi dari aksi monad (atau komonad) di sisi
lain. Teorema Beck~\ref{prop:Beck} memberikan karakterisasi yang praktis bagi
monadisitas.

% segment-id: o014.li2.chapter7.monoidal-algebras.p008
Langkah berikutnya ialah teori Morita dalam \S\ref{sec:Morita}, suatu tema
klasik dalam teori gelanggang. Bagian pertama membahas funktor antara
kategori modul kiri atau kanan yang mempertahankan $\varinjlim$ kecil (atau
$\varprojlim$). Teorema~\ref{prop:Morita} menyatakan bahwa semuanya berasal
dari pengambilan hasil kali tensor (atau $\Hom$) dengan suatu bimodul $P$.
Hal ini sekaligus menegaskan kedudukan istimewa bimodul dalam teori gelanggang
dan modul: bimodul menyediakan sarana yang cukup untuk ``mengubah gelanggang''.
Dari sini diperoleh karakterisasi teoretis-modul bagi ekuivalensi kategori
modul kiri $A\dcate{Mod}$ dan $B\dcate{Mod}$, yang disebut ekuivalensi Morita
(Definisi--Proposisi~\ref{def:Morita-equiv}); karakterisasi untuk modul kiri
dan kanan sepenuhnya sama.

% segment-id: o014.li2.chapter7.monoidal-algebras.p009
Untuk pasangan adjungsi yang ditentukan oleh bimodul-$(A,B)$ $P$, pada bagian
kedua kita menelaah monad dan komonad terkait beserta modul-modulnya. Dengan
asumsi bahwa $P$ merupakan modul-$B$ projektif yang dibangkitkan secara hingga,
Proposisi~\ref{prop:Morita-comonad} menyatakannya secara eksplisit sebagai
aljabar dalam kategori monoidal $(A,A)\dcate{Mod}$ dan koaljabar dalam
$(B,B)\dcate{Mod}$. Dari landasan ini mudah diperoleh penyempurnaan ekuivalensi
Morita (Teorema~\ref{prop:Morita-refined},
Korolari~\ref{prop:Morita-refined-symm}).

% segment-id: o014.li2.chapter7.monoidal-algebras.p010
Teori Morita membahas kategori modul yang diberikan. Sebagai pelengkap, dalam
\S\ref{sec:module-cat} kita mengkaji cara agar suatu kategori abelian
$\mathcal{A}$ ekuivalen dengan kategori modul kanan atas suatu gelanggang $R$.
Jawabannya melibatkan pembangkit projektif kompak dari $\mathcal{A}$
(Teorema~\ref{prop:identify-ModR}). Kita juga akan memberikan syarat cukup
untuk versi yang dibangkitkan secara hingga
(Teorema~\ref{prop:identify-ModR-fg}).

% segment-id: o014.li2.chapter7.monoidal-algebras.p011
Kembali ke teorema monadisitas Beck, paruh pertama
\S\ref{sec:descent-modules} memberikan penerapan sederhana dalam teori modul,
yang disebut penurunan datar untuk modul. Untuk homomorfisme gelanggang
komutatif $K\to L$, hasil ini menjelaskan cara merekonstruksi modul-$K$ dari
modul-$L$ yang dikenai aksi komonad perubahan gelanggang $\mathbf{L}$. Teknik
penurunan banyak digunakan dalam geometri, dan paruh kedua
\S\ref{sec:descent-modules} merumuskan kembali syarat aksi $\mathbf{L}$ dalam
bentuk yang lazim digunakan dalam geometri. Semua ini bertumpu pada persiapan
dalam \S\ref{sec:Morita}.

% segment-id: o014.li2.chapter7.monoidal-algebras.p012
Sebagai kasus khusus penurunan datar, dalam \S\ref{sec:descent-Galois} diambil
$L|K$ sebagai perluasan Galois medan, lalu hasil sebelumnya dirumuskan kembali
melalui aksi grup Galois $\Gal(L|K)$. Teknik ini disebut penurunan Galois.
Teorema utamanya, Teorema~\ref{prop:Galois-descent}, menyatakan bahwa kategori
ruang vektor-$K$ ekuivalen dengan kategori ruang vektor-$L$ yang dilengkapi
aksi semilinear mulus $\Gal(L|K)$. Latihan akan memberikan bukti langsung
lainnya.

% segment-id: o014.li2.chapter7.monoidal-algebras.p013
Berdasarkan penurunan datar atau kasus khususnya, penurunan Galois, berbagai
struktur seperti aljabar dan modul dapat diturunkan lebih lanjut. Dalam
\S\ref{sec:Galois-H1} kita meninjau sejenis data yang dinyatakan melalui ruang
vektor beserta pangkat tensornya. Kemudian, dengan teknik penurunan Galois dan
kohomologi grup, klasifikasi data tersebut direduksi menjadi kohomologi grup
nonabelian $\Hm^1$. Walaupun tujuannya hanya memperagakan teknik, hal ini sudah
membantu menjernihkan beberapa persoalan dasar aljabar, misalnya klasifikasi
bentuk kuadratik takdegenerat.

% segment-id: o014.li2.chapter7.monoidal-algebras.rn001
\begin{petunjukbacaan}
	Bab ini secara garis besar terbagi menjadi tiga bagian. Bagian pertama
(\S\S\ref{sec:alg-in-monoidal-cat}---\ref{sec:bialgebra}) memperkenalkan aljabar
	di atas kategori monoidal, struktur-dg, dan aljabar Hopf; tema-tema ini
	memiliki tempat yang alami dan penting dalam aljabar. Bagian kedua
	(\S\S\ref{sec:Beck}---\ref{sec:module-cat}) berpusat pada monad dan teori
	Morita terkait. Bagian ketiga
	(\S\S\ref{sec:descent-modules}---\ref{sec:Galois-H1}) membahas metode
	penurunan. Ketiganya kurang lebih disusun berurutan, tetapi struktur-dg dan
	kategori-dg tidak diperlukan untuk pembahasan berikutnya. \CHref{sec:simplicial}
	dan \CHref{sec:duality} dalam buku ini akan menggunakan beberapa hasil dari
	bagian kedua. Bahan lanjutan lainnya terdapat dalam latihan bab ini.
\end{petunjukbacaan}

% segment-id: o014.li2.chapter7.monoidal-algebras.h002
\section{Aljabar di Atas Kategori Monoidal}\label{sec:alg-in-monoidal-cat}

% segment-id: o014.li2.chapter7.monoidal-algebras.p014
Untuk suatu gelanggang komutatif $\Bbbk$, kita mengetahui bahwa aljabar-$\Bbbk$
merupakan struktur perkalian yang ditumpangkan pada modul-$\Bbbk$. Definisinya
dapat dinyatakan menggunakan bifunktor hasil kali tensor
$\otimes=\otimes_{\Bbbk}$ pada $\Bbbk\dcate{Mod}$ dan diagram komutatif.
Intinya ialah $(\Bbbk\dcate{Mod},\otimes)$ membentuk kategori monoidal dengan
unit $\munit:=\Bbbk$. Ketika membahas aljabar-$\Bbbk$ komutatif dan berbagai
variannya, peranan sentral dipegang oleh struktur anyaman
$c(X,Y):X\otimes Y\rightiso Y\otimes X$ pada kategori monoidal ini, dengan
sifat simetris $c(Y,X)c(X,Y)=\identity$. Lihat \cite[\S 7.4]{Li1}.

% segment-id: o014.li2.chapter7.monoidal-algebras.p015
\index{kategori monoidal (monoidal category)}
Setelah dirumuskan dalam bahasa yang sesuai, teori ini mencakup
\emph{kategori monoidal} umum dan \emph{kategori monoidal simetris}. Mengikuti
\cite[Definisi 3.1.1, 3.1.2]{Li1}, konvensi bagian ini\footnote{Definisi dalam
literatur lain mungkin sedikit berbeda, tetapi konsep kategori monoidal yang
dihasilkan tetap ekuivalen.} mendefinisikan kategori monoidal sebagai data
$(\mathcal{V},\otimes,a,\munit,\iota)$, yang sering disingkat menjadi
$\mathcal{V}$ atau $(\mathcal{V},\otimes)$, dengan
% segment-id: o014.li2.chapter7.monoidal-algebras.l001
\begin{compactitem}
	\item $\otimes: \mathcal{V} \times \mathcal{V} \to \mathcal{V}$ adalah bifunktor;
	\item $a(X, Y, Z): (X \otimes Y) \otimes Z \rightiso X \otimes (Y \otimes Z)$ adalah morfisme kanonik yang disebut kendala asosiativitas ($X, Y, Z \in \Obj(\mathcal{V})$), dan memenuhi aksioma segilima;
	\item $\munit = \munit_{\mathcal{V}} \in \Obj(\mathcal{V})$ adalah unit;
	\item $\iota: \munit \otimes \munit \rightiso \munit$.
\end{compactitem}
Dari sini diperoleh keluarga isomorfisme kanonik
$\munit \otimes X \xrightarrow[\sim]{\lambda_X} X \xleftarrow[\sim]{\rho_X} X \otimes \munit$;
menurut \cite[Lema 3.1.5]{Li1}, isomorfisme-isomorfisme ini memenuhi
$\lambda_{\munit}=\iota=\rho_{\munit}$. Semuanya disebut kendala unit dan
selanjutnya tidak akan dituliskan secara eksplisit.

% segment-id: o014.li2.chapter7.monoidal-algebras.p016
Suatu \emph{struktur anyaman} pada kategori monoidal adalah isomorfisme antara
bifunktor, yang dituliskan sebagai
% segment-id: o014.li2.chapter7.monoidal-algebras.q001
\[ c = \left( c(X, Y): X \otimes Y \rightiso Y \otimes X\right)_{X, Y \in \Obj(\mathcal{V})}, \]
yang disyaratkan kompatibel dengan kendala asosiativitas dan unit; lihat
\cite[Definisi 3.3.1]{Li1}. Kategori monoidal yang dilengkapi struktur anyaman
tertentu disebut \emph{kategori monoidal beranyam}. Jika
$c(Y,X)c(X,Y)=\identity_{X\otimes Y}$ selalu berlaku, data tersebut disebut
\emph{kategori monoidal simetris}.
\index{struktur anyaman (braid structure)}
\index[sym1]{cXY@$c(X, Y)$}
\index{kategori monoidal!beranyam (braided)}
\index{kategori monoidal!simetris (symmetric)}

% segment-id: o014.li2.chapter7.monoidal-algebras.d001
\begin{definisi}[Aljabar]\label{def:algebra-monoidal}
	\index{aljabar (algebra)}
	Aljabar di atas kategori monoidal $\mathcal{V}$ berarti data
	$(A,\mu,\eta)$, dengan
	% segment-id: o014.li2.chapter7.monoidal-algebras.q002
	\[ A \in \Obj(\mathcal{V}), \quad \mu = \mu_A: A \otimes A \to A, \quad \eta = \eta_A: \munit \to A, \]
	sedemikian sehingga diagram-diagram berikut komutatif:
	% segment-id: o014.li2.chapter7.monoidal-algebras.g003
	\begin{equation*}\begin{gathered}
		\begin{tikzcd}
			\munit \otimes A \arrow[r, "{\eta \otimes \identity}"] \arrow[rd, "\sim"' sloped] & A \otimes A \arrow[d, "\mu"] & A \otimes \munit \arrow[l, "{\identity \otimes \eta}"'] \arrow[ld, "\sim"' sloped] \\
			& A &
		\end{tikzcd} \quad
		\begin{tikzcd}[column sep=tiny]
			(A \otimes A) \otimes A \arrow[rr, "{a(A, A, A)}", "\sim"'] \arrow[d, "{\mu \otimes \identity}"'] & & A \otimes (A \otimes A) \arrow[d, "{\identity \otimes \mu}"] \\
			A \otimes A \arrow[r, "\mu"'] & A & A \otimes A \arrow[l, "\mu"] .
		\end{tikzcd}
	\end{gathered}\end{equation*}
	Dengan demikian, $\mu$ dapat dipandang sebagai operasi perkalian pada $A$,
	sedangkan $\eta$ sebagai morfisme unitnya.
	
	Morfisme dari aljabar $(A,\mu,\eta)$ ke $(A',\mu',\eta')$ didefinisikan
	sebagai morfisme $\phi:A\to A'$ dalam $\mathcal{V}$ yang membuat diagram
	berikut komutatif:
	% segment-id: o014.li2.chapter7.monoidal-algebras.g004
	\[\begin{tikzcd}
		\munit \arrow[r, "\eta"] \arrow[rd, "{\eta'}"'] & A \arrow[d, "\phi"] \\
		& A'
	\end{tikzcd} \quad \begin{tikzcd}
		A \otimes A \arrow[d, "\mu"'] \arrow[r, "{\phi \otimes \phi}"] & A' \otimes A' \arrow[d, "{\mu'}"] \\
		A \arrow[r, "\phi"'] & A' .
	\end{tikzcd}\]
\end{definisi}

% segment-id: o014.li2.chapter7.monoidal-algebras.p017
Kita sering menyingkat $(A,\mu,\eta)$ menjadi $A$. Dalam beberapa literatur,
aljabar juga disebut gelanggang atau monoid di dalam $\mathcal{V}$, dan
terkadang disebut aljabar asosiatif untuk menekankan hukum asosiativitas yang
ditampilkan oleh diagram komutatif kedua.

% segment-id: o014.li2.chapter7.monoidal-algebras.p018
Konsep modul memiliki generalisasi yang serupa.

% segment-id: o014.li2.chapter7.monoidal-algebras.d002
\begin{definisi}[Modul]\label{def:module-algebra}
	\index{modul (module)}
	Misalkan $(A,\mu,\eta)$ aljabar di atas $\mathcal{V}$. Modul-$A$ kiri
	didefinisikan sebagai data $(M,\mu_M)$, dengan
	% segment-id: o014.li2.chapter7.monoidal-algebras.q003
	\[ M \in \Obj(\mathcal{V}), \quad \mu_M: A \otimes M \to M, \]
	yang membuat diagram-diagram berikut komutatif:
	% segment-id: o014.li2.chapter7.monoidal-algebras.g005
	\[\begin{tikzcd}
		\munit \otimes M \arrow[r, "{\eta \otimes \identity}"] \arrow[rd, "\sim"' sloped] & A \otimes M \arrow[d, "{\mu_M}"] \\
		& M
	\end{tikzcd} \quad \begin{tikzcd}[column sep=small]
		(A \otimes A) \otimes M \arrow[rr, "{a(A, A, M)}", "\sim"'] \arrow[d, "{\mu_M \otimes \identity}"'] & & A \otimes (A \otimes M) \arrow[d, "{\identity \otimes \mu_M}"] \\
		A \otimes M \arrow[r, "{\mu_M}"'] & M & A \otimes M \arrow[l, "{\mu_M}"] .
	\end{tikzcd}\]
	Jadi $\mu_M$ dapat dipandang sebagai perkalian skalar pada modul.
	
	Sesuai konvensi, modul $(M,\mu_M)$ disingkat menjadi $M$. Morfisme dari $M$
	ke $M'$ didefinisikan sebagai morfisme $\psi:M\to M'$ dalam $\mathcal{V}$
	yang membuat diagram berikut komutatif:
	% segment-id: o014.li2.chapter7.monoidal-algebras.g006
	\[\begin{tikzcd}
		A \otimes M \arrow[d, "{\mu_M}"'] \arrow[r, "{\identity \otimes \psi}"] & A \otimes M' \arrow[d, "{\mu_{M'}}"] \\
		M \arrow[r, "\psi"'] & M' .
	\end{tikzcd}\]
	
	Modul-$A$ kanan didefinisikan dengan cara serupa. Jika $M$ sekaligus memiliki
	struktur modul-$A$ kiri dan modul-$B$ kanan yang membuat diagram berikut
	komutatif,
	% segment-id: o014.li2.chapter7.monoidal-algebras.g007
	\[\begin{tikzcd}[column sep=small]
		(A \otimes M) \otimes B \arrow[rr, "{a(A, M, B)}", "\sim"'] \arrow[d, "{\mu_M^A \otimes \identity}"'] & & A \otimes (M \otimes B) \arrow[d, "{\identity \otimes \mu_M^B}"] \\
		M \otimes B \arrow[r, "{\mu_M^B}"'] & M & A \otimes M , \arrow[l, "{\mu_M^A}"]
	\end{tikzcd}\]
	maka $M$ disebut bimodul-$(A,B)$; morfisme antara dua bimodul semacam ini
	ialah morfisme $\psi:M\to M'$ yang sekaligus merupakan morfisme modul-$A$
	kiri dan modul-$B$ kanan.
\end{definisi}

% segment-id: o014.li2.chapter7.monoidal-algebras.p019
Jika $M$ merupakan modul-$A$ kiri dan $N$ modul-$B$ kanan, maka $M\otimes N$
secara alami menjadi bimodul-$(A,B)$.

% segment-id: o014.li2.chapter7.monoidal-algebras.r001
\begin{catatan}[Unit sebagai aljabar]\label{rem:unit-algebra}
	Sebagai contoh, $\munit$ menjadi aljabar dengan
	$\mu_{\munit}:=\iota$ dan $\eta_{\munit}:=\identity_{\munit}$. Diagram
	komutatif yang diperlukan mengikuti sifat-sifat standar $\munit$ dan
	$\iota:\munit\otimes\munit\rightiso\munit$; lihat
	\cite[Lema 3.1.5]{Li1}. Untuk setiap aljabar $A$ terdapat tepat satu
	morfisme $\munit\to A$, dan morfisme itu haruslah $\eta_A$.
	
	Serupa dengan itu, mudah dilihat bahwa setiap $M\in\Obj(\mathcal{V})$
	memiliki tepat satu struktur modul-$\munit$ kiri (atau kanan), yaitu dengan
	mengambil $\mu_M:\munit\otimes M\to M$ sebagai kendala unit dari
	$\mathcal{V}$.
\end{catatan}

% segment-id: o014.li2.chapter7.monoidal-algebras.p020
Tujuan berikutnya ialah membahas
% segment-id: o014.li2.chapter7.monoidal-algebras.l002
\begin{inparaenum}[(a)]
	\item komutativitas aljabar,
	\item hasil kali tensornya.
\end{inparaenum}
Keduanya berkaitan erat dan sama-sama memerlukan struktur anyaman pada
$\mathcal{V}$. Kita mulai dari komutativitas.

% segment-id: o014.li2.chapter7.monoidal-algebras.d003
\begin{definisi}\label{def:opposite-braided-algebra}
	\index[sym1]{Aop}
	Dengan menggunakan berbagai sifat alami struktur anyaman, terlihat bahwa
	jika $A$ merupakan aljabar di atas kategori monoidal beranyam $\mathcal{V}$,
	maka mengganti $\mu$ dengan $\mu\circ c(A,A)$ dan membiarkan data lain tetap
	memberikan lagi aljabar di atas $\mathcal{V}$. Aljabar ini didefinisikan
	sebagai \emph{aljabar lawan} $A^{\opp}$ dari $A$. Verifikasinya panjang
	tetapi tidak sulit, dan diserahkan sebagai latihan bab ini.
\end{definisi}

% segment-id: o014.li2.chapter7.monoidal-algebras.p021
Dari funktorialitas struktur anyaman, jika $\phi:A_1\to A_2$ merupakan
morfisme, maka ia juga merupakan morfisme dari $A_1^{\opp}$ ke
$A_2^{\opp}$.

% segment-id: o014.li2.chapter7.monoidal-algebras.d004
\begin{definisi}\label{def:calg-monoidal}
	\index{aljabar!komutatif (commutative)}
	Untuk aljabar $A$ di atas kategori monoidal beranyam $\mathcal{V}$, $A$
	disebut \emph{aljabar komutatif} jika diagram berikut komutatif dalam
	$\mathcal{V}$:
	% segment-id: o014.li2.chapter7.monoidal-algebras.g008
	\[\begin{tikzcd}
		A \otimes A \arrow[rd, "{\mu_A}"'] \arrow[rr, "{c(A, A)}", "\sim"'] & & A \otimes A . \arrow[ld, "{\mu_A}"] \\
		& A &
	\end{tikzcd}\]
	Dengan kata lain, kita mensyaratkan perkalian $\mu_A$ memenuhi hukum
	komutatif. Rumusan ekuivalennya adalah $A=A^{\opp}$.
\end{definisi}

% segment-id: o014.li2.chapter7.monoidal-algebras.p022
Aljabar $\munit$ dari Catatan~\ref{rem:unit-algebra} merupakan contoh sederhana
aljabar komutatif; diagram komutatif yang diperlukan mengikuti
kompatibilitas struktur anyaman $c$ dengan kendala unit.

% segment-id: o014.li2.chapter7.monoidal-algebras.p023
Tetap misalkan $\mathcal{V}$ kategori monoidal beranyam. Seperti dalam kasus
gelanggang komutatif, untuk dua aljabar sembarang $A$ dan $B$, struktur anyaman
$c$ memberikan struktur aljabar kanonik pada $A\otimes B$:
% segment-id: o014.li2.chapter7.monoidal-algebras.l003
\begin{itemize}
	\item Perkalian $\mu_{A \otimes B}$ diambil sebagai komposisi
	% segment-id: o014.li2.chapter7.monoidal-algebras.g009
	\begin{multline*}
		(A \otimes B) \otimes (A \otimes B) \rightiso (A \otimes (B \otimes A)) \otimes B \\
		\xrightarrow{(\identity_A \otimes c(B, A)) \otimes \identity_B} (A \otimes (A \otimes B)) \otimes B \\
		\rightiso (A \otimes A) \otimes (B \otimes B) \xrightarrow{\mu_A \otimes \mu_B} A \otimes B ,
	\end{multline*}
	semua panah tanpa label berasal dari kendala asosiativitas. Dalam notasi
	diagram anyaman, komposisi ini tergambar sebagai berikut:
	% segment-id: o014.li2.chapter7.monoidal-algebras.g010
	\begin{center}\begin{tikzpicture}[baseline=(braid)]
		\pic[
			braid/number of strands=4,
			braid/every strand/.style={black},
			name prefix=braid
		]
		at (0, 0) {braid = {s_2}};
		\coordinate (braid) at ($(braid-1-s)!0.5!(braid-1-e)$);
		\node[above] at (braid-1-s) {$A$};
		\node[above] at (braid-2-s) {$B$};
		\node[above] at (braid-3-s) {$A$};
		\node[above] at (braid-4-s) {$B$};
		
		\coordinate (MA) at ($(braid-1-e)!0.5!(braid-3-e)$);
		\draw[->] (MA) -- ($(MA) + (0, -2em)$) node[below] {$A$};
		\draw[fill=white] ([yshift=0.5em, xshift=-0.5em] braid-1-e) rectangle ([yshift=-1em, xshift=0.5em] braid-3-e);
		\node at ([yshift=-0.25em] MA) {$\mu_A$};
		
		\coordinate (MB) at ($(braid-2-e)!0.5!(braid-4-e)$);
		\draw[->] (MB) -- ($(MB) + (0, -2em)$) node[below] {$B$};
		\draw[fill=white] ([yshift=0.5em, xshift=-0.5em] braid-2-e) rectangle ([yshift=-1em, xshift=0.5em] braid-4-e);
		\node at ([yshift=-0.25em] MB) {$\mu_B$};
	\end{tikzpicture}\end{center}
	
	\item Unit $\eta_{A \otimes B}$ diambil sebagai komposisi
	% segment-id: o014.li2.chapter7.monoidal-algebras.q004
	\[ \munit \xrightarrow[\sim]{\iota^{-1}} \munit \otimes \munit \xrightarrow{\eta_A \otimes \eta_B} A \otimes B. \]
\end{itemize}

% segment-id: o014.li2.chapter7.monoidal-algebras.p024
Untuk memverifikasi bahwa ini memang aljabar, syarat-syarat dalam
Definisi~\ref{def:algebra-monoidal} dapat diturunkan dari sifat-sifat yang
bersesuaian pada $(A,\mu_A,\eta_A)$ dan $(B,\mu_B,\eta_B)$ hanya dengan
menggunakan $c$ (untuk menukar urutan $\otimes$), $\iota$ (untuk menggandakan
unit), dan kendala asosiativitas (untuk mengubah letak tanda kurung).

% segment-id: o014.li2.chapter7.monoidal-algebras.p025
Jika $M$ merupakan modul-$A$ kiri dan $N$ modul-$B$ kiri, maka $M\otimes N$
secara alami menjadi modul-$(A\otimes B)$ kiri. Konstruksi dan verifikasinya
sama dengan paragraf sebelumnya, dengan $c$ digunakan untuk menukar posisi.
Kasus modul kanan pun serupa.

% segment-id: o014.li2.chapter7.monoidal-algebras.p026
Untuk aljabar $\munit$ yang didefinisikan dalam
Catatan~\ref{rem:unit-algebra}, jelas bahwa
$\munit\otimes A\simeq A\simeq A\otimes\munit$, dengan isomorfisme yang
diberikan oleh kendala unit $\mathcal{V}$.

% segment-id: o014.li2.chapter7.monoidal-algebras.r002
\begin{catatan}\label{rem:Alg-otimes-functorial}
	Konstruksi di atas funktorial terhadap $A$ dan $B$: jika $A\to A'$ dan
	$B\to B'$ merupakan morfisme aljabar, maka demikian pula morfisme terkait
	$A\otimes B\to A'\otimes B'$. Fakta lain yang panjang tetapi tidak sulit
	ialah bahwa bila $A,B,C$ semuanya aljabar, kendala asosiativitas
	$(A\otimes B)\otimes C\rightiso A\otimes(B\otimes C)$ merupakan morfisme
	aljabar.
\end{catatan}

% segment-id: o014.li2.chapter7.monoidal-algebras.le001
\begin{lema}\label{prop:CAlg-mu}
	Misalkan $A$ aljabar di atas kategori monoidal beranyam $\mathcal{V}$. Maka
	$A$ komutatif jika dan hanya jika perkalian
	$\mu_A:A\otimes A\to A$ merupakan morfisme aljabar.
\end{lema}
% segment-id: o014.li2.chapter7.monoidal-algebras.pf001
\begin{bukti}
	Agar $\mu_A$ merupakan morfisme aljabar, diperlukan dua syarat. Syarat
	mengenai unit ekuivalen dengan komutativitas diagram
	% segment-id: o014.li2.chapter7.monoidal-algebras.g011
	\[\begin{tikzcd}[column sep=large]
		\munit \otimes \munit \arrow[d, "\iota"'] \arrow[r, "{\eta_A \otimes \eta_A}"] & A \otimes A \arrow[d, "{\mu_A}"] \\
		\munit \arrow[r, "{\eta_A}"'] & A
	\end{tikzcd}\]
	dan syarat ini berlaku otomatis karena $A$ adalah aljabar. Syarat yang
	tersisa ialah komutativitas diagram
	% segment-id: o014.li2.chapter7.monoidal-algebras.g012
	\begin{equation*}
		\begin{tikzcd}[column sep=large]
			(A \otimes A) \otimes (A \otimes A) \arrow[r, "{\mu_A \otimes \mu_A}"] \arrow[d, "{\mu_{A \otimes A}}"'] & A \otimes A \arrow[d, "{\mu_A}"] \\
			A \otimes A \arrow[r, "{\mu_A}"'] & A
		\end{tikzcd}
	\end{equation*}
	Identifikasikan sudut kiri atas dengan
	$A\otimes(A\otimes A)\otimes A$ melalui kendala asosiativitas. Berdasarkan
	asosiativitas $\mu_A$, komposisi dalam diagram sepanjang
	\begin{tikzpicture}[scale=0.5, baseline=(O)]
		\draw[->] (0, 0) -- (1, 0) -- (1, -0.7);
		\coordinate (O) at (0, -0.5);
	\end{tikzpicture}
	sama dengan
	% segment-id: o014.li2.chapter7.monoidal-algebras.q005
	\[ A \otimes (A \otimes A) \otimes A \xrightarrow{\identity \otimes \mu_A \otimes \identity} A \otimes A \otimes A \xrightarrow{\text{perkalian}} A . \]
	Di sisi lain, definisi $\mu_{A\otimes A}$ dan asosiativitas $\mu_A$
	menunjukkan bahwa komposisi dalam diagram sepanjang
	\begin{tikzpicture}[scale=0.5, baseline=(O)]
		\draw[->] (0, 0) -- (0, -0.7) -- (1, -0.7);
		\coordinate (O) at (0, -0.5);
	\end{tikzpicture}
	sama dengan
	% segment-id: o014.li2.chapter7.monoidal-algebras.q006
	\[ A \otimes (A \otimes A) \otimes A \xrightarrow{\identity \otimes (\mu_A c(A,A) ) \otimes \identity} A \otimes A \otimes A \xrightarrow{\text{perkalian}} A . \]
	Jadi $\mu_Ac(A,A)=\mu_A$ mengakibatkan diagram semula komutatif, yakni
	$\mu_A$ merupakan morfisme aljabar.
	
	Untuk implikasi sebaliknya, ambil
	$f\in\left\{\mu_A,\mu_Ac(A,A)\right\}$. Kita mempunyai diagram komutatif
	sederhana
	% segment-id: o014.li2.chapter7.monoidal-algebras.g013
	\[\begin{tikzcd}
		A \otimes (A \otimes A) \otimes A \arrow[r, "{\identity \otimes f \otimes \identity}"] & A \otimes A \otimes A \arrow[r, "{\text{perkalian}}"] & A . \\
		A \otimes A \arrow[u, "{\eta_A \otimes \identity \otimes \eta_A}"] \arrow[r, "f"'] & A \arrow[ru, "\identity"'] \arrow[u, "{\eta_A \otimes \identity \otimes \eta_A}"] &
	\end{tikzcd}\]
	Jika $\mu_A$ merupakan morfisme aljabar, pembahasan sebelumnya menunjukkan
	bahwa baris pertama memberikan komposisi yang sama untuk kedua pilihan $f$.
	Dengan memperhatikan baris kedua diperoleh
	$\mu_A=\mu_Ac(A,A)$.
\end{bukti}

% segment-id: o014.li2.chapter7.monoidal-algebras.le002
\begin{lema}\label{prop:c-as-isom-alg}
	Misalkan $A$ dan $B$ aljabar di atas kategori monoidal simetris
	$\mathcal{V}$. Maka $c(A,B):A\otimes B\to B\otimes A$ merupakan isomorfisme
	aljabar.
\end{lema}
% segment-id: o014.li2.chapter7.monoidal-algebras.pf002
\begin{bukti}
	Dua diagram komutatif harus diverifikasi untuk $c(A,B)$. Diagram pertama
	ialah
	% segment-id: o014.li2.chapter7.monoidal-algebras.g014
	\[\begin{tikzcd}[row sep=tiny]
		& \munit \otimes \munit \arrow[ld, "\sim" sloped] \arrow[dd, "{c(\munit, \munit)}"] \arrow[r, "{\eta_A \otimes \eta_B}"] & A \otimes B \arrow[dd, "{c(A, B)}"] \\
		\munit & & \\
		& \munit \otimes \munit \arrow[lu, "\sim"' sloped] \arrow[r, "{\eta_B \otimes \eta_A}"'] & B \otimes A
	\end{tikzcd}\]
	Komutativitas bagian kanan merupakan akibat funktorialitas $c$, sedangkan
	bagian kiri mengikuti kompatibilitas $c$ dengan kendala unit
	\cite[(3.12)]{Li1}. Untuk diagram kedua, yang menyatakan bahwa $c(A,B)$
	mempertahankan perkalian, argumennya didasarkan pada simetri struktur anyaman
	dan persis sama dengan bagian akhir \cite[\S 7.4]{Li1}.
\end{bukti}

% segment-id: o014.li2.chapter7.monoidal-algebras.pr001
\begin{proposisi}
	Misalkan $A$ dan $B$ aljabar komutatif di atas kategori monoidal simetris
	$\mathcal{V}$. Maka $A\otimes B$ juga komutatif.
\end{proposisi}
% segment-id: o014.li2.chapter7.monoidal-algebras.pf003
\begin{bukti}
	Cukup ditunjukkan bahwa perkalian $\mu_{A\otimes B}$ merupakan morfisme
	aljabar. Ingat definisi $\mu_{A\otimes B}$ dan terapkan
	Lema~\ref{prop:c-as-isom-alg} untuk menangani $c(B,A)$, asumsi komutativitas
	$A,B$ bersama Lema~\ref{prop:CAlg-mu} untuk menangani $\mu_A,\mu_B$, serta
	Catatan~\ref{rem:Alg-otimes-functorial} tentang funktorialitas $\otimes$.
	Maka setiap bagian komposisi tersebut merupakan morfisme aljabar.
\end{bukti}

% segment-id: o014.li2.chapter7.monoidal-algebras.p027
Di atas telah didefinisikan aljabar, modul, dan morfisme di antaranya. Semua
struktur ini membentuk kategori, dengan notasi
% segment-id: o014.li2.chapter7.monoidal-algebras.t001
\begin{center}\begin{tabular}{|c|c|c|c|c|c|} \hline
	struktur & aljabar & aljabar komutatif & modul-$A$ kiri & modul-$B$ kanan & bimodul-$(A,B)$ \\ \hline
	kategori & $\cate{Alg}(\mathcal{V})$ & $\cate{CAlg}(\mathcal{V})$ & $A\dcate{Mod}$ & $\cated{Mod}B$ & $(A,B)\dcate{Mod}$ \\ \hline
	syarat & {\small $\mathcal{V}$: kategori monoidal} & {\small $\mathcal{V}$: kategori monoidal beranyam} & {\small $A$: aljabar} & {\small $B$: aljabar} & {\small $A, B$: aljabar} \\ \hline
\end{tabular}\end{center}
\index[sym1]{Alg@$\cate{Alg}(\cdot)$}
\index[sym1]{CAlg@$\cate{CAlg}(\cdot)$}
\index[sym1]{Mod}

% segment-id: o014.li2.chapter7.monoidal-algebras.p028
Pembahasan di atas juga menunjukkan bahwa
$\cate{Alg}(\mathcal{V})$ dan $\cate{CAlg}(\mathcal{V})$ pada umumnya
memiliki struktur yang lebih sedikit daripada $\mathcal{V}$. Jika tingkat struktur diberi label
angka, situasinya dapat digambarkan secara kasar sebagai
% segment-id: o014.li2.chapter7.monoidal-algebras.g015
\[\begin{tikzcd}[row sep=small]
	0 & 1 & 2 & \infty \\
	\text{kategori} & \text{kategori monoidal} \arrow[l, "{\cate{Alg}}"'] & \text{kategori monoidal beranyam} \arrow[shift right, l, "{\cate{Alg}}"'] \arrow[shift left, bend left, ll, "{\cate{CAlg}}"'] & \text{kategori monoidal simetris} \arrow[loop right, "{\cate{Alg}, \cate{CAlg}}"]
\end{tikzcd}\]
Selain itu, apabila $\cate{Alg}(\mathcal{V})$ atau
$\cate{CAlg}(\mathcal{V})$ menjadi kategori monoidal terhadap $\otimes$,
unitnya selalu $\munit$.

% segment-id: o014.li2.chapter7.monoidal-algebras.p029
Ingat bahwa Lema~\ref{prop:CAlg-mu} menyatakan bahwa aljabar komutatif di atas
kategori monoidal beranyam $\mathcal{V}$ secara otomatis memberikan aljabar di
atas kategori monoidal $\cate{Alg}(\mathcal{V})$, dengan mengambil perkalian
$A\otimes A\to A$ sama dengan perkalian $A$. Bahkan, inilah satu-satunya
pilihan yang mungkin. Dari sini muncul pengamatan sederhana, menarik, dan
penting: terdapat ekuivalensi kategori
% segment-id: o014.li2.chapter7.monoidal-algebras.e001
\begin{equation}\label{eqn:AlgAlg-CAlg}
	\cate{Alg}\left(\cate{Alg}(\mathcal{V}) \right) \simeq \cate{CAlg}(\mathcal{V}), \quad \mathcal{V}:\;\text{kategori monoidal beranyam}.
\end{equation}
Fakta-fakta ini bertumpu pada apa yang disebut \emph{argumen
Eckmann--Hilton}; pembuktiannya sengaja diserahkan sebagai latihan bab ini
yang dilengkapi petunjuk.

% segment-id: o014.li2.chapter7.monoidal-algebras.e002
\begin{contoh}\label{eg:Cartesius-monoidal}
	Misalkan kategori $\mathcal{C}$ mempunyai hasil kali hingga. Khususnya,
	terdapat hasil kali kosong, yakni objek terminal, yang dinotasikan dengan
	$\munit$. Untuk setiap dua objek $X,Y$ terdapat isomorfisme kanonik
	$c(X,Y):X\times Y\rightiso Y\times X$. Dengan data ini,
	$(\mathcal{C},\times)$ menjadi kategori monoidal simetris dengan unit
	$\munit$. Kita dapat membahas aljabar di atas $\mathcal{C}$, yang juga
	disebut objek monoid di dalam $\mathcal{C}$\footnote{Untuk kategori semacam
	ini, \cite[\S 4.11]{Li1} juga membahas objek grup, tetapi sifat invers tidak
	dapat dirumuskan dalam kategori monoidal umum.}, beserta versi komutatifnya.
	\begin{itemize}
		\item Jika $\mathcal{C}=\cate{Set}$, himpunan satu titik adalah unit $\munit$. Aljabar yang bersesuaian adalah monoid, dan aljabar komutatif yang bersesuaian adalah monoid komutatif.
		\item Jika $\mathcal{C}=\cate{Top}$, aljabar yang bersesuaian adalah monoid topologis, dan aljabar komutatif yang bersesuaian adalah monoid topologis komutatif.
	\end{itemize}
	
	Di sisi lain, ambil gelanggang komutatif $\Bbbk$, beri
	$\Bbbk\dcate{Mod}$ struktur kategori monoidal simetris melalui
	$\otimes=\otimes_{\Bbbk}$, dengan $\Bbbk$ sebagai unit $\munit$ dan
	struktur anyaman yang ditentukan oleh $x\otimes y\mapsto y\otimes x$.
	Aljabar (atau aljabar komutatif) yang bersesuaian tepat merupakan aljabar
	(atau aljabar komutatif) dalam pengertian aljabar biasa; lihat
	\cite[Definisi 7.1.1]{Li1}.
\end{contoh}

% segment-id: o014.li2.chapter7.monoidal-algebras.p030
Dalam situasi konkret ini, perkalian aljabar $A$ sering langsung ditulis
sebagai perkalian unsur $xy=x\cdot y$, dan unit langsung diidentifikasikan
dengan unsur $1_A$, sehingga notasi $\mu_A,\eta_A$ dapat dihilangkan.

% segment-id: o014.li2.chapter7.monoidal-algebras.e003
\begin{contoh}[Modul bergradasi dan aljabar bergradasi]\label{eg:graded-module}
	\index{modul bergradasi, aljabar bergradasi (graded module, graded algebra)}
	Tetapkan kategori abelian $\mathcal{A}$ dan himpunan tak kosong $I$. Objek
	$\mathcal{A}^I$ disebut objek bergradasi-$I$ dan ditulis dalam bentuk
	$M=(M^i)_{i\in I}$, sedangkan morfisme di antaranya ditulis
	$(f^i:M^i\to N^i)_{i\in I}$. Lazimnya $M^i\in\Obj(\mathcal{A})$ (atau
	$f^i\in\Mor(\mathcal{A})$) disebut komponen derajat $i$ dari $M$ (atau
	$f$). Untuk kasus khusus $I=\Z$ atau $\Z^n$, hal ini telah berulang kali
	dijelaskan dalam \S\ref{sec:additive-cplx}, \S\ref{sec:triangular-def},
	\S\ref{sec:grading-filtration}, dan lain-lain.
	
	Demi kekonkretan, untuk sementara ambil $\mathcal{A}:=\Bbbk\dcate{Mod}$,
	dengan $\Bbbk$ gelanggang komutatif. Objek bergradasi-$I$ dalam
	$\mathcal{A}$ juga disebut modul bergradasi-$I$. Misalkan selanjutnya $I$
	monoid. Definisikan bifunktor $\otimes$ pada kategori $\mathcal{A}^I$ dengan
	% segment-id: o014.li2.chapter7.monoidal-algebras.e004
	\begin{equation}\label{eqn:graded-module}
		(M \otimes N)^i := \bigoplus_{jk=i} M^j \otimes N^k, \quad (f \otimes g)^i := \bigoplus_{jk=i} f^j \otimes g^k .
	\end{equation}
	Selain itu, definisikan $\munit$ sebagai $\Bbbk$ yang terkonsentrasi pada
	komponen $i=1_I$ (yakni unit $I$). Data ini menjadikan $\mathcal{A}^I$
	kategori monoidal. Syarat unit $\eta_A:\munit=\Bbbk\to A$ secara otomatis
	memastikan $1_A:=\eta_A(1)\in A^0$.
	
	Dalam kasus khusus ketika $I$ komutatif, kita kembali pada modul
	bergradasi-$I$ dan aljabar bergradasi-$I$ yang diperkenalkan dalam
	\cite[Definisi 7.4.1]{Li1}.

	Tetap ambil $\mathcal{A}=\Bbbk\dcate{Mod}$ dan syaratkan $(I,+)$ sebagai
	monoid komutatif. Setiap homomorfisme aditif
	$\epsilon:I\to\Z/2\Z$ menjadikan $\Bbbk\dcate{Mod}^I$ kategori monoidal
	simetris dengan \emph{struktur anyaman Koszul} yang dicirikan oleh
	% segment-id: o014.li2.chapter7.monoidal-algebras.e005
	\begin{equation}\label{eqn:Koszul-braiding}\begin{aligned}
		M^a \otimes N^b & \to N^b \otimes M^a \\
		x \otimes y & \mapsto (-1)^{\epsilon(a)\epsilon(b)} y \otimes x ,
	\end{aligned}\end{equation}
	dengan $a,b\in I$. Aljabar komutatif yang bersesuaian tepat merupakan
	aljabar $\epsilon$-komutatif dalam pengertian
	\cite[Definisi 7.4.3]{Li1}. Berbagai persamaan atau definisi yang dihasilkan
	darinya secara kolektif disebut \emph{aturan tanda Koszul}. Setiap kali
	posisi tensor ditukar, tandanya harus diubah menurut aturan ini.
	\index{aturan tanda Koszul (Koszul sign rule)}
	
	Sebagai contoh, syarat agar $A$ menjadi aljabar komutatif ditulis secara
	eksplisit sebagai
	% segment-id: o014.li2.chapter7.monoidal-algebras.q007
	\[ x y = (-1)^{\epsilon(a) \epsilon(b)} yx, \quad x \in A^a, \; y \in A^b. \]
	
	Situasi tipikal ialah
	% segment-id: o014.li2.chapter7.monoidal-algebras.q008
	\[ I \in \{\Z, \Z_{\geq 0} \}, \quad \epsilon(a) := a \;\bmod 2 , \]
	dengan aljabar $\epsilon$-komutatif yang bersesuaian disebut aljabar
	bergradasi antikomutatif dalam buku yang dirujuk di atas. Karena struktur
	anyaman ini muncul di mana-mana dalam kajian objek bergradasi-$\Z$, istilah
	\emph{komutatif bergradasi} lebih tepat.
	
	Sekarang kembali ke kategori abelian $\mathcal{A}$ yang lebih umum dan
	$(I,+,\epsilon)$ seperti di atas. Untuk menggeneralisasi konstruksi tadi,
	kita perlu memperluas $\mathcal{A}$ menjadi kategori monoidal simetris
	$(\mathcal{A},\otimes)$ sedemikian sehingga
	% segment-id: o014.li2.chapter7.monoidal-algebras.l004
	\begin{itemize}
		\item $\mathcal{A}$ adalah kategori abelian dan koproduk (yakni jumlah langsung) berbentuk seperti dalam \eqref{eqn:graded-module} ada;
		\item $\otimes:\mathcal{A}\times\mathcal{A}\to\mathcal{A}$ merupakan funktor aditif pada setiap variabel dan mempertahankan koproduk tersebut.
	\end{itemize}
	Dalam keadaan ini, $\otimes$ dan struktur anyaman Koszul pada
	$\mathcal{A}^I$ masih dapat didefinisikan dengan cara yang sama, sehingga
	$(\mathcal{A}^I,\otimes)$ menjadi kategori monoidal simetris. Unitnya tetap
	diberikan oleh unit $\munit$ dari $\mathcal{A}$, dipandang sebagai objek
	bergradasi-$I$ yang terkonsentrasi pada derajat nol.
	
	Untuk $(\mathcal{A},\otimes)$ umum, semua verifikasi sama dengan kasus
	$\Bbbk\dcate{Mod}$. Satu-satunya perbedaan ialah bahwa ``unsur'' suatu objek
	tidak lagi bermakna dalam kasus umum.
\end{contoh}

% segment-id: o014.li2.chapter7.monoidal-algebras.cv001
\begin{konvensi}
	Dalam kasus khusus $\mathcal{A}=\Bbbk\dcate{Mod}$, konvensi umum ialah
	memandang modul bergradasi-$I$ $(M^i)_{i\in I}$ sebagai modul yang dilengkapi
	dekomposisi jumlah langsung $M=\bigoplus_{i\in I}M^i$, sedangkan morfisme di
	antaranya disyaratkan mempertahankan komponen jumlah langsung.
\end{konvensi}

% segment-id: o014.li2.chapter7.monoidal-algebras.e006
\begin{contoh}\label{eg:superspace}
	\index{ruang vektor super (super vector space)}
	Dalam konstruksi di atas ($\mathcal{A}=\Bbbk\dcate{Mod}$), ambil $\Bbbk$
	medan dengan $\mathrm{char}(\Bbbk)\neq2$,
	$I=\Z/2\Z=\{0,1\}$, dan $\epsilon=\identity_{\Z/2\Z}$. Kategori monoidal
	simetris yang bersesuaian dinotasikan
	$\cate{Vect}^-_{\Z/2\Z}(\Bbbk)$; objeknya dapat dipandang sebagai ruang
	vektor-$\Bbbk$ yang dilengkapi dekomposisi jumlah langsung
	$V=V_0\oplus V_1$, dan disebut pula \emph{ruang vektor super}. Dari sinilah
	muncul istilah-istilah yang lazim dalam aljabar linear dan fisika matematis,
	seperti aljabar super dan aljabar super komutatif.
\end{contoh}

% segment-id: o014.li2.chapter7.monoidal-algebras.p031
Terakhir, kita membahas pengaruh funktor monoidal longgar terhadap aljabar.

% segment-id: o014.li2.chapter7.monoidal-algebras.d005
\begin{definisi}\label{def:lax-monoidal}
	\index{funktor!monoidal, monoidal longgar, monoidal longgar kanan (monoidal, lax monoidal, right lax monoidal)}
	Misalkan $\mathcal{V}$ dan $\mathcal{V}'$ kategori monoidal. Menurut
	\cite[Catatan 3.1.8]{Li1}, funktor monoidal longgar kanan dari $\mathcal{V}$
	ke $\mathcal{V}'$, yang di sini disingkat menjadi \emph{funktor monoidal
	longgar}, berarti data $(F,\xi_F,\varphi_F)$, dengan
	$F:\mathcal{V}\to\mathcal{V}$ suatu funktor dan morfisme
	% segment-id: o014.li2.chapter7.monoidal-algebras.q009
	\[ \xi_F: F(\cdot) \otimes F(\cdot) \to F(\cdot \otimes \cdot), \quad \varphi_F: \munit_{\mathcal{V}'} \to F\left( \munit_{\mathcal{V}} \right), \]
	yang tunduk pada serangkaian syarat kompatibilitas dengan struktur monoidal.
	Data $(F,\xi_F,\varphi_F)$ sering disingkat menjadi $F$.
	
	Jika $\mathcal{V}$ dan $\mathcal{V}'$ keduanya kategori monoidal beranyam
	(atau simetris), dan $F$ juga membuat diagram berikut komutatif, maka $F$
	disebut funktor beranyam, atau dikatakan kompatibel dengan struktur anyaman:
	% segment-id: o014.li2.chapter7.monoidal-algebras.g016
	\[\begin{tikzcd}[column sep=huge]
		F(X) \otimes F(Y) \arrow[r, "{\xi_F(X, Y)}"] \arrow[d, "{c'(FX, FY)}"'] & F(X \otimes Y) \arrow[d, "{Fc(X, Y)}"] \\
		F(Y) \otimes F(X) \arrow[r, "{\xi_F(Y, X)}"'] & F(Y \otimes X) .
	\end{tikzcd}\]

	Jika $\xi_F$ dan $\varphi_F$ dalam data funktor monoidal longgar keduanya
	isomorfisme, funktor tersebut disebut \emph{funktor monoidal}.
\end{definisi}

% segment-id: o014.li2.chapter7.monoidal-algebras.p032
Tentu kita juga dapat membahas morfisme antara funktor-funktor monoidal
longgar, yang disebut pula transformasi alami:
% segment-id: o014.li2.chapter7.monoidal-algebras.q010
\[ \theta = (\theta_X)_{X \in \Obj(\mathcal{V})}: F \to G, \quad F, G: \mathcal{V} \to \mathcal{V}' . \]
Di samping syarat-syarat transformasi alami, kita mensyaratkan bahwa $\theta$
membuat diagram berikut komutatif:
% segment-id: o014.li2.chapter7.monoidal-algebras.g017
\[\begin{tikzcd}[column sep=huge]
	F(X) \otimes F(Y) \arrow[r, "{\xi_F(X, Y)}"] \arrow[d, "{\theta_X \otimes \theta_Y}"'] & F(X \otimes Y) \arrow[d, "{\theta_{X \otimes Y}}"] \\
	G(X) \otimes G(Y) \arrow[r, "{\xi_G(X, Y)}"'] & G(X \otimes Y)
\end{tikzcd} \quad \begin{tikzcd}
	\munit_{\mathcal{V}'} \arrow[r, "{\varphi_F}"] \arrow[rd, "{\varphi_G}"'] & F(\munit_{\mathcal{V}}) \arrow[d, "{\theta_{\munit_{\mathcal{V}}}}"] \\
	& G(\munit_{\mathcal{V}}) .
\end{tikzcd}\]

Jika $F$ dan $G$ merupakan funktor monoidal, definisi di atas sebenarnya
ekuivalen dengan \cite[Definisi 3.1.10]{Li1}; di sana $\varphi_F$ dan
$\varphi_G$ disebut isomorfisme ``standar''. Dengan demikian dapat dibahas
ekuivalensi antara kategori monoidal atau kategori monoidal beranyam.

% segment-id: o014.li2.chapter7.monoidal-algebras.pr002
\begin{proposisi}\label{prop:lax-monoidal-Alg}
	Funktor monoidal longgar $F$ menginduksi funktor
	$\cate{Alg}(\mathcal{V})\to\cate{Alg}(\mathcal{V}')$. Jika $A$ aljabar di
	atas $\mathcal{V}$, maka $F$ menginduksi
	$A\dcate{Mod}\to F(A)\dcate{Mod}$, dan seterusnya.

	Jika selain itu $\mathcal{V}$ dan $\mathcal{V}'$ kategori monoidal beranyam
	dan $F$ funktor beranyam, maka $F$ juga menginduksi
	$\cate{CAlg}(\mathcal{V})\to\cate{CAlg}(\mathcal{V}')$, sedangkan
	$\cate{Alg}(\mathcal{V})\to\cate{Alg}(\mathcal{V}')$ merupakan funktor
	monoidal longgar. Jika lebih lanjut $\mathcal{V},\mathcal{V}'$ keduanya
	simetris, maka $\cate{CAlg}(\mathcal{V})\to\cate{CAlg}(\mathcal{V}')$ juga
	demikian.
\end{proposisi}
% segment-id: o014.li2.chapter7.monoidal-algebras.pf004
\begin{bukti}
	Misalkan $A$ aljabar di atas $\mathcal{V}$. Definisikan perkalian
	$\mu_{FA}$ sebagai
	$FA\otimes FA\to F(A\otimes A)\xrightarrow{F\mu_A}FA$, dan unit $\eta_{FA}$
	sebagai
	$\munit_{\mathcal{V}'}\to F(\munit_{\mathcal{V}})\xrightarrow{F\eta_A}FA$.
	Verifikasi selebihnya bersifat rutin.
\end{bukti}

% segment-id: o014.li2.chapter7.monoidal-algebras.e007
\begin{contoh}\label{eg:ev0-monoidal}
	Tinjau kategori abelian monoidal simetris $\mathcal{A}$ dan $\mathcal{A}^I$
	dari Contoh~\ref{eg:graded-module}. Definisikan funktor
	% segment-id: o014.li2.chapter7.monoidal-algebras.q011
	\[ \mathrm{ev}_0: \mathcal{A}^I \to \mathcal{A}, \quad \left(M^i\right)_{i \in I} \mapsto M^0. \]
	Funktor ini eksak sekaligus monoidal longgar simetris: cukup ambil
	% segment-id: o014.li2.chapter7.monoidal-algebras.q012
	\[ \xi_{\mathrm{ev}_0}(M, N): M^0 \otimes N^0 \to (M \otimes N)^0 := \bigoplus_{i+j=0} M^i \otimes N^j \]
	sebagai morfisme yang jelas, dan
	$\varphi_{\mathrm{ev}_0}:\munit\to\munit$ sebagai identitas. Dengan demikian,
	$A\mapsto A^0$ menginduksi funktor monoidal longgar
	% segment-id: o014.li2.chapter7.monoidal-algebras.q013
	\[ \cate{Alg}(\mathcal{A}^I) \to \cate{Alg}(\mathcal{A}), \quad \cate{CAlg}(\mathcal{A}^I) \to \cate{CAlg}(\mathcal{A}). \]
\end{contoh}

% segment-id: o014.li2.chapter7.monoidal-algebras.p033
Catatan berikut membahas aljabar di atas kategori monoidal umum $\mathcal{V}$.
Ingat bahwa ordinal hingga berkorespondensi satu-satu dengan $\Z_{\geq0}$;
ordinal hingga $\mathbf{n}$ yang bersesuaian dengan $n$ diidentifikasikan
dengan himpunan terurut total berunsur $n$, yaitu
$\{0,\ldots,n-1\}$ dengan $0\leq1\leq2\leq\cdots$.

% segment-id: o014.li2.chapter7.monoidal-algebras.r003
\begin{catatan}[Aljabar berjalan]\label{rem:walking-algebra}
	\index{aljabar!berjalan (walking)}
	\index[sym1]{FinOrd@$\cate{FinOrd}$}
	Pemetaan $f:A\to B$ antara himpunan terurut sebagian disebut
	mempertahankan urutan jika
	$a\leq a'\implies f(a)\leq f(a')$. Selanjutnya, tuliskan $\cate{FinOrd}$
	untuk kategori ordinal hingga yang morfismenya adalah pemetaan
	mempertahankan urutan antara ordinal. Kategori ini menjadi kategori
	monoidal ketat terhadap
	$\mathbf{n}\otimes\mathbf{m}:=\mathbf{n}+\mathbf{m}$. Morfisme
	$f\otimes g:\mathbf{n}\otimes\mathbf{m}\to\mathbf{n}'\otimes\mathbf{m}'$
	memetakan $n$ unsur pertama (atau $m$ unsur terakhir) dengan $f$ (atau $g$),
	dan objek $\mathbf{0}$ adalah unit\footnote{Di sini $\mathbf{1}$ selalu
	menyatakan ordinal; unit kategori monoidal $\mathcal{V}$ dinotasikan terpisah
	dengan $\munit_{\mathcal{V}}$.}. Jadikan $\mathbf{1}$ aljabar di atas
	$\cate{FinOrd}$: ambil unit $\emptyset=\mathbf{0}\hookrightarrow\mathbf{1}$
	dan perkalian
	$\mathbf{1}\otimes\mathbf{1}=\mathbf{2}\twoheadrightarrow\mathbf{1}$.
	Diagram komutatif yang terkait sama sekali tidak sulit diverifikasi.
	
	Ordinal $\mathbf{1}$ (atau $\mathbf{2}$) dapat dipandang sebagai objek
	(atau morfisme) ``berjalan''. Alasannya, untuk sembarang kategori
	$\mathcal{C}$, menentukan funktor $\mathbf{1}\to\mathcal{C}$ (atau
	$\mathbf{2}\to\mathcal{C}$) setara dengan menentukan objek (atau morfisme)
	dalam $\mathcal{C}$. Sekarang tinjau kategori monoidal $\mathcal{V}$. Kita
	klaim terdapat isomorfisme kategori
	% segment-id: o014.li2.chapter7.monoidal-algebras.q014
	\[
		\left\{ \mathcal{V}\;\text{-aljabar}\; A \right\} \simeq \left\{ \text{funktor monoidal}\; F: \cate{FinOrd} \to \mathcal{V} \right\},
	\]
	yang juga berarti bahwa $\cate{FinOrd}$ adalah ``aljabar berjalan''.
	Pertama, bagi funktor monoidal
	$F:\cate{FinOrd}\to\mathcal{V}$, definisikan $A:=F(\mathbf{1})$. Sebagai
	kasus khusus sederhana dari Proposisi~\ref{prop:lax-monoidal-Alg}, $A$
	secara alami menjadi aljabar di atas $\mathcal{V}$: perkalian
	$\mu:A\otimes A\to A$ adalah $F(\mathbf{2}\twoheadrightarrow\mathbf{1})$,
	sedangkan unit $\eta:\munit_{\mathcal{V}}\to A$ adalah
	$F(\mathbf{0}\hookrightarrow\mathbf{1})$.
	
	Sebaliknya, dari aljabar $(A,\mu,\eta)$ di atas $\mathcal{V}$, definisikan
	funktor monoidal $F:\cate{FinOrd}\to\mathcal{V}$ dengan
	% segment-id: o014.li2.chapter7.monoidal-algebras.q015
	\[ F(\mathbf{n}) = A^{\otimes n} := \underbracket{A \otimes (A \otimes (\cdots))}_{n\;\text{salinan}\; A}, \quad A^{\otimes 0} := \munit_{\mathcal{V}}, \]
	Untuk $f:\mathbf{n}\to\mathbf{m}$, jika $f$ injektif, maka
	$F(f):A^{\otimes n}\to A^{\otimes m}$ diperoleh dengan menyisipkan
	$\eta:\munit_{\mathcal{V}}\to A$ di luar citra $f$; jika $f$ surjektif,
	$F(f)$ diperoleh dengan memakai $\mu:A\otimes A\to A$ untuk menggabungkan
	secara bertahap faktor-faktor dalam setiap serat. Kendala asosiativitas dan
	sifat-sifat lain kategori monoidal memastikan operasi ini terdefinisi dengan
	baik. Pemetaan umum $f$ memiliki faktorisasi surjektif--injektif yang unik,
	dan dengan cara ini $F$ didefinisikan pada morfisme. Funktorialitasnya
	jelas. Tidak sulit memverifikasi bahwa kedua konstruksi saling invers.
\end{catatan}
